\chapter{Connection to the cuspidal thesis problem}\label{ch:cusps}

The preceding chapters prove statements about the two specified
nodal moduli problems. They also identify the work needed for a
new moduli problem of maps whose source curves may have ordinary
cusps. They do not construct such a moduli stack by replacing one
word in a definition.

\section{What transfers formally, and what needs geometry}

An ordinary cusp over an algebraically closed field of characteristic
zero has completed local ring
\[
                   k[[x,y]]/(y^2-x^3)=k[[t^2,t^3]].
\]
Unlike a node, it has one branch on normalization. The normalization
parameterization is $x=t^2,y=t^3$. This local model specifies the
allowed singularity; it does not specify a stability condition for
an entire pointed curve or map.

To begin the thesis construction, first write down the objects:
proper flat finitely presented source families, allowed geometric
fibers, distinct smooth labelled sections, a map to the fixed target,
the meaning of $\beta$, and an explicit stability condition. We
must know whether an elliptic subcurve, a contracted cuspidal
component, or a rational component with too few attachments is
allowed. None of those decisions is determined by the word ``cusp.''

Once these data are chosen, the earlier work separates into the
following tasks.

\begin{enumerate}
\item \emph{Formal category theory.}
Use pairs $(u,h)$ with cartesian source squares, preservation of
markings, and $f=g h$. Identities and composition have the same
explicit formulas as before. The projection retains the parameter
scheme. Its fibers are groupoids because the arrows over an identity
are required to be isomorphisms. These checks need no new theorem
about cusps.

\item \emph{Base change of the proposed objects.}
The cartesian-arrow calculation transfers once pulling back an object
really stays inside the proposed class. Properness, flatness, and
finite presentation already behave correctly. One must verify the
geometric singularity and stability conditions over arbitrary
geometric fibers and arbitrary base changes. A condition stated
only for closed complex points is not yet enough to define our
category over all schemes.

\item \emph{The sheaf condition for isomorphisms.}
The underlying morphism theorem is unchanged. Glue $h_i$, glue
$h_i^{-1}$, and check the equations for the base, sections, and
target map after the cover. Once objects are well defined, this
part does not depend on a node having two branches. It also does
not establish that the isomorphism sheaf is representable or
unramified.

\item \emph{Effective descent.}
Here there is a real geometric input to supply. Nodes and ordinary
cusps are hypersurface singularities, hence Gorenstein. For a proper
flat finitely presented family with these geometric fibers, the
relative-dualizing theorem used earlier therefore still supplies an
invertible $\omega_{C/S}$ and its base-change identifications.
Smooth markings still give relative Cartier divisors. The
candidate bundle
\[
       \omega_{C/S}(\textstyle\sum\sigma_a)\otimes f^*A^{\otimes q}
\]
is consequently functorial for any fixed integer $q$.
What needs a proof for the \emph{chosen} stability condition is its
relative ampleness, or some other effective-descent mechanism.
The nodal formula counting two node branches cannot simply be
copied to a cusp. If a suitable functorial ample bundle is found,
the same polarized descent theorem constructs the total space;
the maps and sections then descend as before. The permitted
singularities and the proposed stability condition must be
recoverable after the chosen covers as well.

\item \emph{Algebraicity and the Deligne--Mumford property.}
These are further theorems. One needs a bounded embedding result
and a representable parameter locus for the permitted cuspidal
families and maps (or another algebraicity argument). One must
prove representability of the diagonal and control infinitesimal
automorphisms. Finiteness of automorphisms, and the role of
characteristic zero, must be justified for the new stability
condition. Canonical ampleness alone no longer gives the nodal
argument; see the example below.

\item \emph{Finite type, separatedness, and properness.}
Finite type requires a uniform bound for the chosen discrete data.
Separatedness requires uniqueness of an extension of an
isomorphism. Properness requires existence of admissible limits
after suitable finite base changes, with the desired stability
condition, and compatibility with the target map. Ordinary stable
reduction proves these statements for the nodal compactification;
it is not a proof of them for a different cuspidal compactification.
One must prove or locate the corresponding theorems for the exact
new moduli problem.
\end{enumerate}

In particular, the geometry needed for object descent is different
from the geometry needed for a finite-type atlas or for a unique
limit. Keeping these obligations separate tells us which claim
has actually been established at any point in the thesis.

\section{A small warning supplied by an actual cusp}

\begin{Ex}[Ample pointed dualizing bundle with infinitely many automorphisms]
Consider the projective cuspidal cubic
\[
                  C=\{y^2z=x^3\}\subset\PP^2_{\CC},
                  \qquad p=[0:1:0].
\]
The point $p$ is smooth and the cusp is $[0:0:1]$.
The normalization is $\PP^1$. At the cusp, the quotient
$\CC[[t]]/\CC[[t^2,t^3]]$ has dimension one, with basis
the class of $t$. The normalization sequence therefore gives
$\chi(\mathcal O_C)=\chi(\mathcal O_{\PP^1})-1=0$, so
the arithmetic genus is one. As a hypersurface, $C$ is
Gorenstein, and its dualizing sheaf is invertible. The general
Gorenstein-curve formula
$\deg\omega_C=2p_a(C)-2$, Stacks Project Lemma 53.8.3,
\MCtag{0C19}, gives degree zero. Thus $\omega_C(p)$ has
degree one and is ample, by the componentwise ampleness
criterion \MCtag{0B5Y}. This calculation does not presume
that the curve is nodal.

Nevertheless, for every $\lambda\in\CC^*$,
\[
       [x:y:z]\longmapsto[\lambda^2x:\lambda^3y:z]
\]
preserves $C$ and fixes $p$. These are infinitely many distinct
automorphisms, forming a positive-dimensional group. They also
give nonzero infinitesimal automorphisms by allowing
$\lambda=1+\epsilon$ with $\epsilon^2=0$.

This object is not asserted to belong to any proposed
pseudostable moduli problem. It shows concretely why the nodal
equivalence between the stability inequalities, canonical
ampleness, and finite automorphisms needs reconsideration when
cusps are admitted. For a constant map $C\to X$, these
automorphisms still preserve the map. For the inclusion
$C\hookrightarrow\PP^2$, an automorphism preserving that
inclusion must be the identity. The target map can change the
automorphism calculation dramatically.
\end{Ex}

\begin{Ej}\label{ex:cusp-roadmap}
Suppose you have proved a functorial ample-bundle theorem for
your proposed cuspidal maps and that all defining conditions
are local for \'etale covers. Which of Fantechi's two conditions
can you now prove? Which further assertions are still required
before calling the result a proper Deligne--Mumford stack?
Use the cuspidal cubic to explain why the second question is
not answered merely by ampleness.
\end{Ej}

The next thesis step is therefore to specify the exact cuspidal
stability condition and test these obligations against it. These
notes establish neither the intended pseudostable-map stack nor
a virtual class for it. The later theories named in the project
brief are deliberately left for subsequent notes.
